%% Chapter 1: Introduction
%% Status: skeleton.  Only the learning goals are fixed; they come from the
%% exam pool.  The sectioning is deliberately not invented in advance --
%% it will follow the source material once that has been read.

\chapter{Introduction}
\label{ch:introduction}

This chapter sets the stage. We fix notation, recall the complexity-theoretic
background the rest of the book assumes, and lay out the four strategies for
dealing with computationally hard problems that structure Parts~III--V. Nothing
here is deep; the point is to agree on a common language.

\begin{goals}
  \goalitem{state what is assumed as background and what is not}
  \goalitem{use $O$, $\Theta$, $\Omega$ and in particular $\Ostar$ notation correctly}
  \goalitem{explain informally what a word RAM is and why we usually compute in that model}
  \goalitem{recall the definitions of \textbf{P}, \textbf{NP}, NP-hardness and NP-completeness}
  \goalitem{name the four standard responses to NP-hardness and say when each is appropriate}
\end{goals}

\todo{This chapter is not written yet. Note the division of labour with
\cref{app:preliminaries}, which is written: the appendix carries the
\emph{background} --- notation, graph vocabulary, the $\Ostar$ convention, the
machine model, and the definitions of \textrm{P}, \textrm{NP} and
polynomial-time reduction --- together with a self-check that tells a reader
which of those they need to read. This chapter should point there for all of
it and spend its own pages on the part the appendix deliberately does not
carry: what makes a problem hard, the four responses to hardness that organise
Parts~III--V, and how the parts of the book fit together.}
